\documentclass[10pt]{article}

\usepackage[utf8]{inputenc}

\usepackage[T1]{fontenc}

\usepackage{amsmath}

\usepackage{amsfonts}

\usepackage{amssymb}

\usepackage[version=4]{mhchem}

\usepackage{stmaryrd}



\begin{document}

тегией. В позициях из $X_{0}$ выбор альтернативы случаен. Это случайное блуждание определяет вероятностное распределение на множестве окончательных позиций. Принимая за выигрыш игрока математич. ожидание его выигрыша на окончательных позициях, получают бескоалиционную игру в нормальной форме.



П. и. наз. и г р о й с пол но й ин ф о р м ац и е й, если информационные множества игроков состоят из одной позиции каждое. В игре с полной информацией существует ситуация равновесия в чистых стратегиях (те о ре м а Це р ме лон a).



Важную роль в П. и. играют смешанные стратегии. Две смешанные стратегии игрока наз. ә к в и в а л е в тны м и, если в ситуациях, отличающихся только этими стратегиями, вероятности каждой окончательной позиции равны. Игрок имеет по л н у ю п а м я т ь, если каждое его информационное множество следует за единственной альтернативой всех предшествующих информационных множеств этого игрока. Наличие полной памяти у игрока означает, что в момент соверпения хода он помнит все информационные множества, в к-рых он находился, и свои выборы в них. С т р aте г и е й п о в д е н и я наз. такая сметіанная стратегия, при к-рой случайные выборы альтернатив информационных множеств стохастически независимы. Каждая смешанная стратегия игрока әквивалентна нек-рой стратегии поведения тогда и только тогда, когда игрок имеет полную память (т е о р е м а H a).



Более общими являются П. и. с бесконечным множеством альтернатив в каждой позиции и с бесконечно продолжающимися партиями, а также игры на грабах.



Лum.: [1] Позиционные игры, М., 1967; [2] Б е р ж К. Общая теория игр нескольких лиц, пер. с франц., М., 1961 ; [3] К у н Г. У., в сб.: Позиционные игры, М., 1967, с. 13-40 [4] T h o m p s o $\mathrm{G}$ G. L., в кн.: Contributions to the theory of games, v. 2, Princeton, 1953, p. $267-77$; [5] B о р о б ь е в Н. Н. “Проблемы кибернетики", 1962 , в. 7, с. 5-20; [6] Б у й К о н г К ы о н г, “Вестник ЛГУ”, 1969, № 1, с. 49-59.



ПОИА РАСІРЕДЕЛЕНИЕ - распределение вероятностей случайной величины $X_{n}$, принимающей целые неотрицательные значения $k, 0 \leqslant k \leqslant n$, в соответствии с формулой



\[

\begin{gathered}

\boldsymbol{P}\left\{X_{n}=k\right\}= \\

=C_{n}^{k} \frac{b(b+c) \ldots[b+(k-1) c] r(r+c) \ldots[r+(n-k-1) c]}{(b+r)(b+r+c) \ldots[b+r+(n-1) c]},

\end{gathered}

\]



где целые $n>0, b>0, r>0, c \geqslant-1$ - параметры, или әквивалентной формулой



\[

P\left\{X_{n}=k\right\}=

\][



=C\_\{n\}\^{}\{k\} \textbackslash frac\{p(p+\textbackslash gamma) ···[p+(k-1) \textbackslash gamma] q(q+\textbackslash gamma) ···[q+(n-k-1) \textbackslash gamma]\}\{(1+\textbackslash gamma) ···[1+(n-1) \textbackslash gamma]\}=



][



=C\_\{(p / \textbackslash psi)+k-1\}\^{}\{k\} C\_\{(q / \textbackslash gamma)+n-k-1\}\^{}\{n-k\} / C\_\{(1 / \textbackslash gamma)+n-1\}\^{}\{n\},

]



где целое $n>0$, действительные $0<p<1, \quad q=1-p$, $\gamma>0$ - параметры. Связь между (1) и (2) устанавливается равенствами



\[

p=\frac{b}{b+r}, \quad q=\frac{r}{b+r}, \quad \gamma=\frac{c}{b+r} .

\]



Математич. ожидание и дисперсия П. р. равны соответственно $\mathrm{E} X_{n}=n p$ и $\mathrm{D} X_{n}=n p q \frac{1+\gamma n}{1+\gamma}$. Сгециальные случаи П. р.: $X_{n}$ имест при $c=0$ биномиальное распределение с пгараметрами $n$ и $p ; X_{n}$ имеет при $s=-1$ гипергеометрическое распределение с параметрами $M=b$, $N=b+r$ и $n$. При $b \rightarrow \infty, r \rightarrow \infty$, когда $p=b /(b+r)$ постоянно, и $\gamma=c(b+r) \rightarrow 0$, П. р. стремится к биномиальному распределению с параметрами $n$ и $p$.



П. р. было рассмотрено Д. Пойа (G. Pólya, 1923) в связи с т. н. урново й схем о й П ой а. Из урны, содержащей $b$ черных и $r$ красных шаров, осуществляется выбор с возвращением при условии, что каждый извлеченный шар возвращается в урну вместе с $c$ шарами того же цвета. Если $X_{n}$ - полное число черных шаров в выборке объема $n$, то распределение $X_{n}$ задается формулами (1) или (2). Последовательность $X_{n}, \quad n=1,2, \ldots$, представляет собой дискретный марковский процесс, причем состояния процесса определяются числом черных шаров в выборке в момент $n$, а условная вероятность перехода от состояния $k$ в момент времени $n$ в состояние $k+1$ в момент времени $n+1$ равна



\[

P\left\{X_{n+1}=k+1 \mid X_{k}=k\right\}=\frac{b+k c}{b+r+n c}=\frac{p+k \psi}{1+n \nu}

\]



(зависит от $n$ ).



Предельным переходом из урновой схемы Пойа может быть получен п р о ц е с с П о й а - Неоднородный марковский процесс с непрерывным временем, принадлежащий классу процессов «чистого размножения». При условии, что за бесконечно малое время $\Delta t$ происходит лишь одно извлөчение шара, при $n \rightarrow \infty$, когда $n p \rightarrow t, n \gamma \rightarrow \alpha t$, выводится предельная условная вероятность перехода из состояния $k$ в состояние $k+1$ за время $\Delta t$ :



$\mathrm{P}\{X(t+\Delta t)=k+1 \mid X(t)=k\}=\frac{1+\alpha k}{1+\alpha t} \Delta t+o(\Delta t)$.



При переходе от урновой схемы Пойа к процессу Пойа возникает важная предельная форма П. р. Именно, вероятность $P_{k}(t)$ в момент времени $t$ пребывать в состоянии $k$ равна



\[

\begin{gathered}

P_{n}(t)=C_{(1 / \alpha)+n-1}^{n}\left(\frac{\alpha t}{1+\alpha t}\right)^{n}\left(\frac{1}{1+\alpha t}\right)^{1 / \alpha} \\

\left(P_{0}(0)=1\right) .

\end{gathered}

\]



Полученное пределььое распределение является отрицательным биномиальным распределением с параметрами $1 / \alpha$ и $1 /(1+\alpha t)$ (соответствующее математич. ожидание равно $t$, а дисперсия $t(1+\alpha t)$.



Урновая модель и процесс Пойа, в к-рых возникает П. р. и его предельная форма, являются моделями с әфф е к т м по с л е д е й с т в и я (извлечение шара определенного цвета из урны увеличивает вероятность извлечения шара того же цвета при следующем испытании).



При стремлении параметра $\alpha$ к нулю процесс Пойа переходит в пуассоновский процесс, а П. р. при $\alpha \rightarrow 0$ имеет своим пределом Пуассона распределение с параметром $t$.



Лum.- [1] Ф е л л е р В., Введение в теорию вероятностей и ее приложения, пер. с англ., 2 изд., т. 1-2, М., 1967.



ПОЙА ТЕОРЕМА: пусть $R^{D}$ множество отображений конечного множества $D,|D|=n$, в множество $R$ и пусть $G$ - групга подстановок множества $D$, порождающая разбиение $R^{D}$ на классы эквивалентности. при к-ром $f_{1}, f_{2} \in R D$ принадлежат одному и тому же классу тогда и только тогда, когда найдется такое $g \in G$, что $f_{1}(g(d))=f_{2}(d)$ для всех $d \in D$; если каждому $r \in R$ сопоставлен вес $w(r)$ - элемент коммутативного кольца (вес $f$ полагается равным $w(f)=\prod_{d \in D} w(f(d)$ ) и вес $w(F)$ класса $F \subset R D$ определяется как вес любого $f \in F)$, то



\[

\begin{gathered}

\sum_{F} w(F)=P\left(G ; \sum_{r \in R} w(r), \sum_{r \in R} w^{2}(r), \ldots\right. \\

\left.\ldots, \sum_{r \in R} w^{n}(r)\right),

\end{gathered}

\]



где в левой части равенства сумма берется по всем классам эквивалентности, а



$P\left(G ; x_{1}, x_{2}, \ldots, x_{n}\right)=|G|^{-1} \sum_{g \in G} \prod_{k=1}^{n} x_{k}^{j(g)}=P(G)$





\end{document}